一个分布的完整信息量太大，没人能拿着密度函数做决策。实际工作里问的通常是三句话：中心在哪里，摊得有多开，两个量是不是一起动。这一章把这三句话对应的数字造出来，并且说清楚每个数字为了简洁而丢掉了什么。

把它们排成一条线索会更好读。期望是分布的重心，也是平方损失下最好的常数预测；方差是用了这个最好的常数之后仍然剩下的误差；协方差是两个变量共同偏离的记录，它的出现是因为方差不像期望那样可以随手相加；条件期望把``常数''升级成``函数''，于是最好的预测随信息而动，方差也随之裂成解释得了和解释不了的两块。

四条判断值得在读之前先摆出来，因为它们最容易被含混地记成别的样子。第一，期望的线性性不需要独立，任何能写成一堆变量之和的计数量都可以逐项求期望，哪怕这些变量彼此纠缠——这是把复杂计数问题变简单的主要手段。第二，方差没有这条便利，\(\Var(X+Y)\) 比两个方差之和多出一个交叉项，那个交叉项就是协方差；而方差的量纲是原量纲的平方，能与数据摆在一起比较的是标准差。第三，相关系数只丈量直线贴合的程度，它等于零允许两个变量之间存在任意强的确定关系，联合 Gaussian 是唯一常用的例外。第四，条件期望是给定信息后平方损失下的最佳预测，几何上就是投影，全方差分解则量出这条信息解释掉了多少变异。

按顺序读四篇是最省力的路径，因为每一篇都在上一篇留下的缺口上开工。若从具体困惑进入，第一篇处理``分布太复杂但只想要平均总量''，第二篇处理``均值一样风险却不一样''以及量纲与标准化，第三篇处理``两个量的关联怎么度量、什么时候度量失灵''，第四篇处理``已经拿到一条信息，预测和误差各该怎么改写''。

\begin{itemize}
\tightlist
\item \hyperref[sec:05-Probability/04-Expectation-and-Moments/01]{``期望、线性性与指示变量''}；
\item \hyperref[sec:05-Probability/04-Expectation-and-Moments/02]{``方差、矩与尺度''}；
\item \hyperref[sec:05-Probability/04-Expectation-and-Moments/03]{``协方差、相关性与条件波动''}；
\item \hyperref[sec:05-Probability/04-Expectation-and-Moments/04]{``条件期望是利用信息的最佳预测''}。
\end{itemize}

读完之后要保留一份警惕：这些数字都是摘要，不是分布。均值与方差相同的两个分布可以有完全不同的尾部；相关系数为零的两个变量可以互相决定；而所有这些量都以相应阶的矩存在为前提，重尾数据上它们可能压根没有定义。每次动用一个摘要时顺手问一句它刻意忽略了什么，比记住公式更重要。
